\SourceLeaf{353}%
schen Theile ihm eigenthümliche in einer Reihe von einzelnen Abhandlungen bekannt zu machen. Es wird dadurch noch der Vortheil gewonnen, dass auf diese Art manche ein selbstständiges Interesse darbietende Untersuchungen, welche mit den übrigen in enger Verwandtschaft stehen, sie vorbereiten und in ein helleres Licht setzen, auch wenn von denselben bei den in Rede stehenden Messungen selbst keine unmittelbare Anwendung gemacht ist, doch mit grösserer Ausführlichkeit entwickelt werden können, als bei dem frühern Plane mit einer gleichmässigen Behandlung der Gegenstände verträglich sein würde.

In die Klasse solcher Untersuchungen gehört namentlich diejenige, welche den Gegenstand der vorliegenden \textit{ersten} Abhandlung ausmacht. Den Hauptinhalt derselben bildet eine Methode, nach welcher ein System von Dreiecken auf der Oberfläche eines Umdrehungs-Ellipsoids, ohne etwas von der Schärfe aufzuopfern, so berechnet werden kann, als wenn es auf einer Kugelfläche sich befände. Diese Methode findet ihre Grundlage in der Auflösung eines viel umfassendern Problems, welche der Verf. in einer 1822 geschriebenen und von Herrn Conferenzrath \textsc{Schumacher} im dritten Heft der Astronomischen Abhandlungen zum Druck beförderten Denkschrift gegeben hat, unter dem Titel: \textit{Allgemeine Auflösung der Aufgabe, die Theile einer gegebenen Fläche auf einer andern gegebenen Fläche so abzubilden, dass die Abbildung dem Abgebildeten in den kleinsten Theilen ähnlich wird.} Der Verf hat diejenigen Darstellungen einer Fläche auf einer andern, welche der angegebenen Bedingung Genüge leisten, zur Abkürzung des Vortrags und weil sie überhaupt als eine sehr reiche Hülfsquelle für die Rechnungen der höhern Geodaesie eine besondere Benennung wohl verdienen, mit dem Namen \textit{conforme Darstellungen} belegt, welches sonst vage Beiwort also hier immer in einer präcis bestimmten Bedeutung zu verstehen ist. \textsc{Mercators} und die stereographische Projection sind bekannte Beispiele conformer Darstellungen der Kugelfläche auf der Ebene.

Es ist kaum nöthig, zu bemerken, dass die Aehnlichkeit in den kleinsten (unendlich kleinen) Theilen wohl unterschieden werden muss von der Ähnlichkeit in allen endlichen Theilen. Die letztere ist nur in speciellen Fällen zu erreichen möglich, wenn nemlich die erste Fläche entweder auf die zweite selbst oder auf eine ihr ähnliche abgewickelt werden kann; im Allgemeinen aber, wo die Conformität \textit{nur} in der Ähnlichkeit der kleinsten Theile besteht, ist das \textit{Vergrösserungsverhältniss}, d.\,i. das Verhältniss, in welchem die auf beiden Flächen einan%
